%% TRANSCRIBED (PDF 92-97 = printed III-1..III-6, plus the tail of 1.2
%% printed on III-7 = PDF 98, above the §2 heading)
%% NOTE: printed pages III-6 and III-7 are Taylor & Francis blank watermark
%% pages in the working PDF (local PDF 97-98 carry only the publisher logo,
%% cf. the parallel note at the head of ch3_s2.tex).  Everything below whose
%% source is printed III-6, or the top of III-7 before §2, was transcribed
%% from the textually identical 1989 Addison-Wesley photo-reprint (public
%% McGill scan, working/serre-mcgill.pdf, pages 98 and 99), after checking
%% that it dovetails verbatim with III-5 of the working copy (PDF 96) and
%% with the opening of ch3_s2.tex.
%% TODO-VERIFY(mc98-99): the III-6/III-7 passages below follow the McGill
%% reprint scan, not the working PDF (whose pages 97-98 are blank).

%% ---------- Chapter III opening (printed page III-1) ----------
In this Chapter, we define what it means for an abelian \textit{l}-adic
representation to be \emph{locally algebraic} and we prove
(cf.~\S\ref{sec:III.2.3}) that such a representation, when rational,
comes from a linear representation of one of the groups $S_m$ of
Chapter~\ref{chap:II}.

When the ground field is a composite of quadratic extensions of $\Q$,
any rational semi-simple \textit{l}-adic representation is \emph{ipso
facto} locally algebraic; this is proved in \S\ref{sec:III.3}, as a
consequence of a result on transcendental numbers due to Siegel and
Lang.

In the local case, an abelian semi-simple representation is locally
algebraic if and only if it has a ``Hodge--Tate decomposition''.  This
fact, due to Tate (Collège de France, 1966), is proved in the Appendix,
together with some complements.

\sechead{1}{The local case}{THE LOCAL CASE}

\subsechead{1.1}{Definitions}

Let $p$ be a prime number and $K$ a finite extension of $\Q_p$; let
$T = \mathrm{R}_{K/\Q_p}(\mathbf{G}_{m/K})$ be the corresponding
algebraic torus over $\Q_p$ (cf.~Weil \cite{weil-adeles}, Chap.~I).

Let $V$ be a finite dimensional $\Q_p$-vector space and denote, as
usual, by $\GL_V$ the corresponding linear group; it is an algebraic
group over $\Q_p$, and $\GL_V(\Q_p) = \Aut(V)$.

Let $\rho\colon \Gal(\bar K/K)^{\mathrm{ab}} \to \Aut(V)$ be an abelian
$p$-adic representation of $K$ in $V$, where
$\Gal(\bar K/K)^{\mathrm{ab}}$ denotes the Galois group of the maximal
abelian extension of $K$.  If
$i\colon K^* \to \Gal(\bar K/K)^{\mathrm{ab}}$ is the canonical
homomorphism of local class field theory (cf.~for instance
Cassels-Fröhlich \cite{cassels-frohlich}, chap.~VI, \S2), we then get a
continuous homomorphism $\rho\circ i$ of $K^* = T(\Q_p)$ into
$\Aut(V)$.

\begin{dfn}
The representation $\rho$ is said to be locally algebraic if there is
an algebraic morphism $r\colon T \to \GL_V$ such that
$\rho\circ i(x) = r(x^{-1})$ for all $x \in K^*$ close enough to $1$.
\end{dfn}

Note that, if $r\colon T \to \GL_V$ satisfies the above condition, it
is unique; this follows from the fact that any non-empty open set of
$K^* = T(\Q_p)$ is Zariski dense in $T$.  We say that $r$ is the
algebraic morphism \emph{associated with} $\rho$.

\begin{examples}
\begin{nlist}
\item Take $K = \Q_p$, and $\dim V = 1$, so that $\rho$ is given by a
continuous homomorphism
$\Gal(\bar\Q_p/\Q_p)^{\mathrm{ab}} \to U_p$, where $U_p$ is the group
of $p$-adic units.  It is easy to see that there exists an element
$\nu \in \Z_p$ such that $\rho\circ i(x) = x^\nu$ if $x$ is close
enough to $1$.

The representation $\rho$ is locally algebraic if and only if $\nu$
belongs to $\Z$.  This happens for instance when $V = V_p(\mu)$,
cf.~Chap.~\ref{chap:I}, \S\ref{sec:I.1.2}, in which case $\nu = -1$ and
$r$ is the canonical one-dimensional representation of
$T = \mathbf{G}_{m/\Q_p}$.

\item The abelian representation associated to a Lubin--Tate formal
group (cf.~\cite{lubin-tate} and \cite{cassels-frohlich}, chap.~VI,
\S3) is locally algebraic (and also of the form $u \mapsto u^{-1}$ on
the inertia group).
\end{nlist}
\end{examples}

\begin{proposition}[1]\label{prop:III.1.1.1}
Let $\rho\colon \Gal(\bar K/K)^{\mathrm{ab}} \to \Aut(V)$ be a locally
algebraic abelian representation of $K$.  The restriction of $\rho$ to
the inertia subgroup of $\Gal(\bar K/K)^{\mathrm{ab}}$ is semi-simple.
\end{proposition}

Let us identify the inertia subgroup of $\Gal(\bar K/K)^{\mathrm{ab}}$
with the group $U_K$ of units of $K$.  By assumption, there is an open
subgroup $U'$ of $U_K$ and an algebraic morphism $r$ of $T$ into
$\GL_V$ such that $\rho(x) = r(x^{-1})$ if $x \in U'$.  Let $W$ be a
sub-vector space of $V$ stable by $\rho(U_K)$; it is then stable by
$\rho(U')$, hence by $r(T)$.  But every linear representation of a
torus is semi-simple.  Hence, there exists a projector
$\pi\colon V \to W$ which commutes with the action of $T$.  If we put
$\pi' = \frac{1}{(U_K:U')}\sum\limits_{s \in U_K/U'}
\rho(s)\,\pi\,\rho(s^{-1})$, we obtain a projector $\pi'\colon V \to W$
which commutes with all $\rho(s)$, $s \in U_K$, q.e.d.

Conversely, let us start from a representation $\rho$ whose restriction
to $U_K$ is semi-simple.  If we make a suitable large finite extension
$E$ of $\Q_p$, the restriction of $\rho$ to $U_K$ may be brought into
diagonal form, i.e. is given by continuous characters
$\chi_i\colon U_K \to E^*$, $i = 1, \ldots, n$.  We assume $E$ large
enough to contain all conjugates of $K$, and we denote by $\Gamma_K$
the set of all $\Q$-embeddings of $K$ into $E$.  Recall
(cf.~chap.~\ref{chap:II}, \S\ref{sec:II.1.1}) that the $[\sigma]$,
$\sigma \in \Gamma_K$, make a basis of the character group $X(T)$ of
$T$.

\begin{proposition}[2]\label{prop:III.1.1.2}
The representation $\rho$ is locally algebraic if and only if there
exist integers $n_\sigma(i)$ such that
\[
\chi_i(u) = \prod_{\sigma \in \Gamma_K} \sigma(u)^{-n_\sigma(i)}
\]
for all $i$ and all $u$ close enough to $1$.
\end{proposition}

The necessity is trivial.  Conversely, if there exist such integers
$n_\sigma(i)$, they define algebraic characters
$r_i = \prod [\sigma]^{n_\sigma(i)}$ of $T$, hence a linear
representation $r$ of $T/E$.  It is clear that there is an open
subgroup $U'$ of $U_K$, such that $\rho(u) = r(u^{-1})$ for all
$u \in U'$.  Hence it remains to see that $r$ can be defined over
$\Q_p$ (cf.~chap.~\ref{chap:II}, \S\ref{sec:II.2.4}).  But the trace
$\theta_r = \sum r_i$ of $r$ (loc.~cit.) is such that
$\theta_r(u) \in \Q_p$ for all $u \in U'$.  Since $U'$ is
Zariski-dense in $T$, this implies that $\theta_r$ is ``defined over
$\Q_p$'', hence that $r$ can be defined over $\Q_p$ (loc.~cit.),
q.e.d.

\runhead{Extension of the ground field}

Let $K'$ be a finite extension of $K$, and let $\rho'$ be the
restriction of the given representation $\rho$ to $\Gal(\bar K/K')$.
Then $\rho'$ is locally algebraic if and only $\rho$ is; moreover, if
this is so, the associated algebraic morphisms
\[
r\colon T \to \GL_V\,, \qquad r'\colon T' \to \GL_V
\]
are such that $r' = N\circ r$, where $T'$ is the torus associated with
$K'$ and $N_{K'/K}\colon T' \to T$ is the algebraic morphism defined by
the norm from $K'$ to $K$.
%% NOTE(p095): "if and only $\rho$ is" is AS PRINTED (the second "if" is
%% missing in the original, in both the 1968 printing and the McGill
%% reprint).
%% TODO-VERIFY(p095): the print shows here an indented line "K'/K" between
%% "the torus associated with" and "K' and N_{K'/K}", in both the 1968
%% printing and the McGill reprint; it appears to be a typewriter
%% typesetting false start duplicating the subscript of $N_{K'/K}$.  The
%% prose above follows the mathematically coherent reading ("the torus
%% associated with K'", cf. the book's own usage "the Q_p-torus associated
%% to K", p. IV-37); the stray "K'/K" is not part of the sentence.

All this follows easily from the commutativity of the diagram
\[
\begin{array}{ccc}
K^* & \longrightarrow & \Gal(\bar K/K)^{\mathrm{ab}}\\
\uparrow N & & \uparrow\\
K'^* & \longrightarrow & \Gal(\bar K/K')^{\mathrm{ab}}\,,
\end{array}
\]
and from the fact that the kernel of $N_{K'/K}\colon T' \to T$ is
connected for the Zariski topology.

\runhead{Exercise}

Give an example of a locally algebraic abelian $p$-adic representation
of dimension $2$ which is not semi-simple.

\subsechead{1.2}{Alternative definition of ``locally algebraic'' via Hodge--Tate modules}

Let us recall first the notion of a \emph{Hodge--Tate module}
(cf~\cite{serre-galpdiv}, \S2); here $K$ is only assumed to be complete
with respect to a discrete valuation, with perfect residue field $k$
and $\mathrm{char}(K) = 0$, $\mathrm{char}(k) = p$.  Denote by $C$ the
\emph{completion} $\hat K$ \emph{of the algebraic closure} of $K$.

The group $G = \Gal(\bar K/K)$ acts continuously on $\bar K$.  This
action extends continuously to $C$.  Let $W$ be a $C$-vector space of
finite dimension upon which $G$ acts continuously and semi-linearly
according to the formula
\[
s(cw) = s(c).s(w) \qquad (s \in G,\ c \in C \text{ and } w \in W).
\]

Let $\chi\colon G \to U_p$ be the homomorphism of $G$ into the group
$U_p = \Z_p^*$ of $p$-adic units, defined by its action on the
$p^\nu$-th roots of unity (cf.~chap.~\ref{chap:I},
\S\ref{sec:I.1.2}):
\[
s(z) = z^{\chi(s)} \qquad \text{if } s \in G \text{ and }
z^{p^\nu} = 1.
\]

Define for every $i \in \Z$ the subspace
\[
W^i = \{\, w \in W \mid sw = \chi(s)^i w \text{ for all } s \in G \,\}
\]
of $W$.  This is a $K$-vector subspace of $W$.  Let
$W(i) = C \tens_K W^i$.  This is a $C$-vectorspace upon which $G$ acts
in a natural way (i.e. by the formula $s(c \tens y) = s(c) \tens
s(y)$).  The inclusion $W^i \to W$ extends uniquely to a $C$-linear map
$\alpha_i\colon W(i) \to W$, which commutes with the action of $G$.

\begin{proposition}[(\textnormal{\textit{Tate}})]
Let $\coprod W(i)$ be the direct sum of the $W(i)$, $i \in \Z$.  Let
$\alpha\colon \coprod W(i) \to W$ be the sum of the $\alpha_i$'s
defined above.  Then $\alpha$ is injective.
\end{proposition}

For the proof see \cite{serre-galpdiv}, \S2, prop.~4.

\begin{corollary}
The $K$-spaces $W^i$ ($i \in \Z$) are of finite dimension.  They are
linearly independent over $C$.
\end{corollary}

\begin{dfn}[1]\label{dfn:III.1.2.1}
The module $W$ is of Hodge--Tate type if the homomorphism
$\alpha\colon \coprod\limits_{i \in \Z} W(i) \to W$ is an isomorphism.
\end{dfn}

Let now $V$ be as in \S\ref{sec:III.1.1}, a vector space over $\Q_p$,
of finite dimension.  Let $\rho\colon G \to \Aut(V)$ be a $p$-adic
representation.  Let $W = C \tens_{\Q_p} V$ and let $G$ act on $W$ by
the formula
\[
s(c \tens v) = s(c) \tens \rho(s)(v), \qquad s \in G,\ c \in C,\
v \in V.
\]

\begin{dfn}[2]
The representation $\rho$ is of Hodge--Tate type if the $C$-space
$W = C \tens_{\Q_p} V$ is of Hodge--Tate type
(cf.~def.~\ref{dfn:III.1.2.1}).
\end{dfn}

\runhead{Example}

Let $F$ be a $p$-divisible group of finite height
(cf.~\cite{serre-pdiv}, \cite{tate-pdiv}); let $T$ be its Tate module
(loc.~cit.) and $V = \Q_p \tens T$.  The group $G$ acts on $V$, and
Tate has proved (\cite{tate-pdiv}, Cor.~2 to Th.~3) that this Galois
module is of Hodge--Tate type; more precisely, one has
$W = W(0) \oplus W(1)$, where $W = C \tens V$ as above.

\begin{theorem}[(\textnormal{\textit{Tate}})]
Assume $K$ is a finite extension of $\Q_p$ (i.e. its residue field is
finite).  Let $\rho\colon G \to \Aut(V)$ be an abelian $p$-adic
representation of $K$.  The following properties are equivalent:
\begin{alist}
\item $\rho$ is locally algebraic (cf.~\S\ref{sec:III.1.1}).

\item $\rho$ is of Hodge--Tate type and its restriction to the inertia
group is semi-simple.
\end{alist}
\end{theorem}

For the proof, see the Appendix.
